% Options for packages loaded elsewhere
% Options for packages loaded elsewhere
\PassOptionsToPackage{unicode}{hyperref}
\PassOptionsToPackage{hyphens}{url}
\PassOptionsToPackage{dvipsnames,svgnames,x11names}{xcolor}
%
\documentclass[
  letterpaper,
  DIV=11,
  numbers=noendperiod]{scrartcl}
\usepackage{xcolor}
\usepackage{amsmath,amssymb}
\setcounter{secnumdepth}{-\maxdimen} % remove section numbering
\usepackage{iftex}
\ifPDFTeX
  \usepackage[T1]{fontenc}
  \usepackage[utf8]{inputenc}
  \usepackage{textcomp} % provide euro and other symbols
\else % if luatex or xetex
  \usepackage{unicode-math} % this also loads fontspec
  \defaultfontfeatures{Scale=MatchLowercase}
  \defaultfontfeatures[\rmfamily]{Ligatures=TeX,Scale=1}
\fi
\usepackage{lmodern}
\ifPDFTeX\else
  % xetex/luatex font selection
\fi
% Use upquote if available, for straight quotes in verbatim environments
\IfFileExists{upquote.sty}{\usepackage{upquote}}{}
\IfFileExists{microtype.sty}{% use microtype if available
  \usepackage[]{microtype}
  \UseMicrotypeSet[protrusion]{basicmath} % disable protrusion for tt fonts
}{}
\makeatletter
\@ifundefined{KOMAClassName}{% if non-KOMA class
  \IfFileExists{parskip.sty}{%
    \usepackage{parskip}
  }{% else
    \setlength{\parindent}{0pt}
    \setlength{\parskip}{6pt plus 2pt minus 1pt}}
}{% if KOMA class
  \KOMAoptions{parskip=half}}
\makeatother
% Make \paragraph and \subparagraph free-standing
\makeatletter
\ifx\paragraph\undefined\else
  \let\oldparagraph\paragraph
  \renewcommand{\paragraph}{
    \@ifstar
      \xxxParagraphStar
      \xxxParagraphNoStar
  }
  \newcommand{\xxxParagraphStar}[1]{\oldparagraph*{#1}\mbox{}}
  \newcommand{\xxxParagraphNoStar}[1]{\oldparagraph{#1}\mbox{}}
\fi
\ifx\subparagraph\undefined\else
  \let\oldsubparagraph\subparagraph
  \renewcommand{\subparagraph}{
    \@ifstar
      \xxxSubParagraphStar
      \xxxSubParagraphNoStar
  }
  \newcommand{\xxxSubParagraphStar}[1]{\oldsubparagraph*{#1}\mbox{}}
  \newcommand{\xxxSubParagraphNoStar}[1]{\oldsubparagraph{#1}\mbox{}}
\fi
\makeatother


\usepackage{longtable,booktabs,array}
\usepackage{calc} % for calculating minipage widths
% Correct order of tables after \paragraph or \subparagraph
\usepackage{etoolbox}
\makeatletter
\patchcmd\longtable{\par}{\if@noskipsec\mbox{}\fi\par}{}{}
\makeatother
% Allow footnotes in longtable head/foot
\IfFileExists{footnotehyper.sty}{\usepackage{footnotehyper}}{\usepackage{footnote}}
\makesavenoteenv{longtable}
\usepackage{graphicx}
\makeatletter
\newsavebox\pandoc@box
\newcommand*\pandocbounded[1]{% scales image to fit in text height/width
  \sbox\pandoc@box{#1}%
  \Gscale@div\@tempa{\textheight}{\dimexpr\ht\pandoc@box+\dp\pandoc@box\relax}%
  \Gscale@div\@tempb{\linewidth}{\wd\pandoc@box}%
  \ifdim\@tempb\p@<\@tempa\p@\let\@tempa\@tempb\fi% select the smaller of both
  \ifdim\@tempa\p@<\p@\scalebox{\@tempa}{\usebox\pandoc@box}%
  \else\usebox{\pandoc@box}%
  \fi%
}
% Set default figure placement to htbp
\def\fps@figure{htbp}
\makeatother





\setlength{\emergencystretch}{3em} % prevent overfull lines

\providecommand{\tightlist}{%
  \setlength{\itemsep}{0pt}\setlength{\parskip}{0pt}}



 


\KOMAoption{captions}{tableheading}
\makeatletter
\@ifpackageloaded{caption}{}{\usepackage{caption}}
\AtBeginDocument{%
\ifdefined\contentsname
  \renewcommand*\contentsname{Table of contents}
\else
  \newcommand\contentsname{Table of contents}
\fi
\ifdefined\listfigurename
  \renewcommand*\listfigurename{List of Figures}
\else
  \newcommand\listfigurename{List of Figures}
\fi
\ifdefined\listtablename
  \renewcommand*\listtablename{List of Tables}
\else
  \newcommand\listtablename{List of Tables}
\fi
\ifdefined\figurename
  \renewcommand*\figurename{Figure}
\else
  \newcommand\figurename{Figure}
\fi
\ifdefined\tablename
  \renewcommand*\tablename{Table}
\else
  \newcommand\tablename{Table}
\fi
}
\@ifpackageloaded{float}{}{\usepackage{float}}
\floatstyle{ruled}
\@ifundefined{c@chapter}{\newfloat{codelisting}{h}{lop}}{\newfloat{codelisting}{h}{lop}[chapter]}
\floatname{codelisting}{Listing}
\newcommand*\listoflistings{\listof{codelisting}{List of Listings}}
\makeatother
\makeatletter
\makeatother
\makeatletter
\@ifpackageloaded{caption}{}{\usepackage{caption}}
\@ifpackageloaded{subcaption}{}{\usepackage{subcaption}}
\makeatother
\usepackage{bookmark}
\IfFileExists{xurl.sty}{\usepackage{xurl}}{} % add URL line breaks if available
\urlstyle{same}
\hypersetup{
  colorlinks=true,
  linkcolor={blue},
  filecolor={Maroon},
  citecolor={Blue},
  urlcolor={Blue},
  pdfcreator={LaTeX via pandoc}}


\author{}
\date{}
\begin{document}


\section{Data Analysis for Business
Applications}\label{data-analysis-for-business-applications}

\section{MA321 Syllabus Section 85B Fall
2026}\label{ma321-syllabus-section-85b-fall-2026}

\subsection{Info}\label{info}

3 Credits\\
Mon 6:30pm - 9:20pm\\
Room: C409

\subsection{Instructor Information}\label{instructor-information}

calvin\_williamson@fitnyc.edu\\
office: B831 Science and Math\\
office hours: M 1-3, W 11-12, R 12-1, or appointment

\subsection{Description}\label{description}

This course covers intermediate statistics topics with applications to
business. Students graph, manipulate, and interpret data using
statistical methods and Excel. Topics include data transformations,
single and multiple regression, time series, analysis of variance, and
chi-square tests. Applications are from the areas of retail, finance,
management, and marketing. Prerequisite(s): MA 222

\subsection{Outcomes}\label{outcomes}

Upon completion of this course, students will be able to:

\begin{enumerate}
\def\labelenumi{\arabic{enumi}.}
\tightlist
\item
  Graphically display data using EXCEL.
\item
  Mathematical manipulation of data using formulas in EXCEL.
\item
  Plotting and analyzing time series graphs.
\item
  Investigating trend, cyclical and seasonal components of time series.
\item
  Applying smoothing techniques for forecasting with time series.
\item
  Correlation and graphing regression line.
\item
  Simple linear regression.
\item
  Multiple linear regression.
\end{enumerate}

\subsection{Course Materials}\label{course-materials}

\subsubsection{Textbook}\label{textbook}

Some readings are from an OER textbook that is free. No other textbook
is required.

\subsubsection{Software}\label{software}

We will be using Google Spreadsheets or other free software for all work
in this course. Since these are web-based applications there is NO OTHER
SOFTWARE required for the course besides a web browser.

\subsection{Topics}\label{topics}

\begin{itemize}
\tightlist
\item
  Simple Regression
\item
  Multiple Regression
\item
  Time Series
\item
  Seasonality
\item
  Applications of Normal Distributions
\item
  Inventory Models
\item
  Economic Order Quantity
\item
  Newsvendor Problem
\item
  Empirical Probability Distributions
\end{itemize}

\subsection{Evaluation}\label{evaluation}

Your grade will come from these parts:

\begin{itemize}
\tightlist
\item
  Quizzes (88\%)
\item
  In Class Work Credits (8\%)
\item
  Homework Credits (4\%)
\end{itemize}

Each of these parts is described in more detail below

\subsubsection{Quizzes}\label{quizzes}

Your quiz grade will come from 5 quizzes, roughly covering 2 or 3 weeks
of material each. The quizzes are 30 minutes each and are usually 5 or 6
questions each. These quizzes are with no notes, no internet, no phone,
no software, no AI tools. Pen and paper and calculator only. They are
some multiple choice, some short answer, some true false.

\subsubsection{In Class Work Credits (8\%) (1-3 per
class)}\label{in-class-work-credits-8-1-3-per-class}

These are credits you obtain for demonstrating you have completed
assigned problems during class.

These assignments are done during class and you show them to me as you
complete them. You will earn a single credit for each successful
assignment completion. You must be in attendance to earn these credits.
I decide if the work is complete enough to receive credit for.

The in class work can only be turned in by showing it to me. There is
not a way to turn these assignments in by emailing them to me or
submitting in Brightspace, or any other digital way of submitting them.

You show them to me and I mark them as complete when you are in class.

You may miss up to 2 classes and you can still show me the assignments
from those classes when you are in class again. But you may not do this
more than 2 weeks worth.

\subsubsection{Homework Credits (4\%)}\label{homework-credits-4}

These are credits you obtain for the problems assigned between classes.

You show me the homework at the beginning of the class and you will earn
a single credit for each successful assignment completion. As with the
in class work, homework can only be turned in by showing it to me, not
by emailing it or submitting it in Brightspace.

Some homeworks will be counted for a grade, but not always.

\subsection{Midterm, Final Exam}\label{midterm-final-exam}

There is NO MIDTERM and NO FINAL EXAM in this course.

\subsection{AI Policy}\label{ai-policy}

Use of AI in this course is forbidden, with one exception. You may use
AI to work with the content from the textbook: to explain a passage,
work through an example, or check your understanding of a topic we are
covering.

Every other use of AI is not allowed. That includes the in class work
and the homework, which are your own work. The quizzes are pen and paper
with no notes, no internet, no phone, no software, and no AI tools.

\subsection{Institutional Policies and
Resources}\label{institutional-policies-and-resources}

The academic honor code, withdrawal procedures, emergency preparedness,
student conduct policy, and academic and technical support resources all
apply to this course:

\href{https://calvinw.github.io/FIT-courses/shared/InstitutionalPoliciesAndResources.html}{Institutional
Policies and Resources}




\end{document}
